\documentclass[10pt]{article}
\usepackage[margin=0.55in]{geometry}
\usepackage{amsmath, amssymb}
\usepackage{enumitem}
\usepackage{titlesec}
\setlist{nosep, leftmargin=1.2em}
\titlespacing*{\section}{0pt}{5pt}{2pt}
\titleformat{\section}{\normalsize\bfseries}{}{0pt}{}
\setlength{\parskip}{3pt}
\setlength{\parindent}{0pt}
\pagestyle{empty}

\begin{document}

\begin{center}{\large\bfseries $0.3\,M_o$ Lower-Bound on Unbalanced Moment}\end{center}

\section*{Why a floor exists}
Any analysis that hands us $M_{sc}$ --- classical EFM, FEA, or a design-strip integration --- can drift low for reasons unrelated to the real slab behavior: an under-resolved mesh, a symmetric-tributary cancellation that ignores pattern loading, or an instability that leaves a plate singularity undamped. For a punching-shear checker this is a safety concern: the eccentric-shear stress scales with $M_{sc}$, so a bogus low $M_{sc}$ shows a column passing when it would fail under a fully-loaded panel. The Direct Design Method in ACI 318 provides a well-accepted lower bound for exactly this case, and we adopt it as an unconditional floor on the per-column per-axis moment used in the check.

\section*{The static moment $M_o$}
Per ACI 318-19 \S 8.10.3.2 the total factored static moment in one span direction for a flat two-way panel is
\[
M_o \;=\; \frac{w_u\, l_2\, l_n^2}{8},
\]
with $w_u$ the factored load per unit area, $l_2$ the perpendicular span (center-to-center of supports), and $l_n$ the clear span in the direction of bending, $l_n \geq 0.65\,l_1$. We compute an $M_o$ in each of the two principal directions (x-bending and y-bending), giving two floors per column.

\section*{What the app does}
For every column, we estimate nearest-neighbor spacings $l_{+x}, l_{-x}, l_{+y}, l_{-y}$ by scanning the column layout. A missing neighbor (edge or corner column) is replaced by the distance from the column to the slab boundary along that direction. The worst-case adjacent panel in each direction is used for $l_1$; the averaged transverse spacing is used for $l_2$. Clear spans are computed as $l_n = \max(0.65\,l_1,\, l_1 - c_i)$ per ACI.

Given $M_o$ in each direction we apply
\[
|M_{u,\text{used}}|
\;=\; \max\!\bigl(\,|M_{u,\text{solver}}|,\, 0.3\, M_o\,\bigr),
\]
per column, per axis, and feed the result into the eccentric-shear stress term in \S 22.6.7. Where the solver result is larger, it governs unchanged; where it is smaller (or zero), the $0.3\,M_o$ floor takes over. Sign is preserved from the solver when present.

\section*{Why $0.3\,M_o$}
The DDM historical provision ($0.3\,M_o$) approximates the unbalanced moment produced at a column by a one-sided live-load pattern on a single panel, assuming $\alpha$ = 2 and $w_L / w_D \approx 1$. It is conservative for uniformly-loaded panels with stiff slabs, and roughly tight for flexible slabs with live-dead ratios near 1. It is not a ceiling --- real Mu from pattern loading plus span asymmetry can be several times $M_o$ --- which is why we keep the solver's value whenever it exceeds the floor.

\section*{What this does not do}
\begin{itemize}
\item Does not replace the solver. Solver Mu always wins when it exceeds $0.3\,M_o$. The floor is a safety net, not a design value.
\item Does not redistribute. $0.3\,M_o$ is applied per axis per column independently; cross-axis coupling is whatever the solver reports.
\item Does not verify the reinforcement capacity. If $0.3\,M_o$ governs, the column-strip flexural design must carry that moment --- check separately per \S 8.4.2.2.2.
\item Does not apply to non-gravity loadings. Seismic unbalanced moments, crane loads, etc., are governed by their own lower bounds and should be entered directly as overrides.
\end{itemize}

\section*{Interaction with Table 8.4.2.2.4}
When the direct-shear gate passes for a column/axis (per Table 8.4.2.2.4 and the accompanying whitepaper), we set $\gamma_v = 0$ on that axis. The eccentric-shear contribution drops out regardless of $M_u$, so the $0.3\,M_o$ floor does not affect the DCR for that axis. Where the gate fails, the floored $M_u$ contributes via $\gamma_v \cdot M_u \cdot c / J_c$.

\end{document}
